\chapter{Weak Disorder, Strong Disorder } \label{ch:3}
\section{The normalized partition function} In this section, we introduce a natural martingale that will play an important  role in many results concerning the asymptotic behavior of the polymer. 

For any fixed path of the brownian motion, $\{\om(V_t)\}_{t\geq 0}$ is a Poisson process of intensity $\nu r^d$ and has associated exponential martingales $\{\exp \lef(\b \om(V_t)-\lambda(\b)\nu r^d t\rig)\}_{t\geq 0}$. Hence, for	$t\geq 0$, the \emph{normalized partition function}
\begin{equation} \label{eq:Wt}
W_t = e^{-\lambda (\b)\nu r^d t} Z_t,
\end{equation}
defines a positive, mean $1$, c\`adl\`ag martingale with respect to $\{\mathcal{G}_t\}_{t\geq 0}$.

By Doob's martingale convergence theorem \cite[Chapter 2, Corollary 2.11]{revuzYor}, we get the existence of a random variable $W_\infty$ such that
\begin{equation}
W_\infty = \lim_{t\to\infty} W_t \  \text{ a.s.}
\end{equation} 

\begin{theorem}
There is a dichotomy: either the limit $W_\infty$ is almost-surely positive, or it is almost-surely zero. Otherwise stated, we have either
\begin{equation}  \label{eq:WD}
\IP \{W_\infty > 0\} = 1,
\end{equation}
or
\begin{equation} \label{eq:SD}
\IP \{W_\infty = 0\} = 1.
\end{equation}
\end{theorem}
\begin{proof}
Denote by $e_t$ the renormalized weight 
\begin{equation}
e_t = \exp(\b \om(V_t(B)) - \lambda(\b)\nu r^d t)
\end{equation}
By the Markovian property \eqref{eq:Markov2}, we get that for all positive times $t$ and $s$,
\begin{equation} \label{eq:WtIntegralEquation}
W_{s+t} = P [e_{t} \, W_s \circ \theta_{t,B_t} ].
\end{equation}
In Section \ref{sec:proofSelfconsistent}, we will justify that one can take the limit as $s\to\infty$ in this equality, in order to get that
\begin{equation} \label{eq:WinfinityIntegralEquation2}
W_\infty = P [e_{t} \, W_\infty \circ \theta_{t,B_t} ].
\end{equation}
Then, notice that \eqref{eq:WinfinityIntegralEquation2} also writes 
\begin{equation} \label{eq:WinfinityIntegralEquation}
W_\infty = W_t \, \int_{\mathbb{R}^d} P_t^{\b,\om} (B_t \in \mathrm{d} x)\, W_\infty \circ \theta_{t,x}.
\end{equation}
Since $W_t > 0$ $\IP$-a.s and since $P_t^{\b,\om}$ has positive density with respect to Lebesgue's measure, we obtain by \eqref{eq:WinfinityIntegralEquation} that
\[ \forall t>0, \quad \{W_\infty = 0\} = \{W_\infty \circ \theta_{t,x} = 0, \ x\text{-a.e.}\}, \]
or, equivalently,
\[  \{W_\infty = 0\} = \left\{ \int_{\mathbb{R}^d} P (B_t \in \mathrm{d} x)\, W_\infty \circ \theta_{t,x} =0 \right\}.
 \]
The event of the right-hand side belong 
%[???] 
% Yes: $\{W_\infty = 0\} \stackrel{a.s.}{=} \{ \int W_\8 \circ \theta_{t,x} P(B_t \in dx) =0\}  \in  \mathcal{G}_{[t,\infty)}$ par discretization de l'$\int$
to the $\sigma$-field $\mathcal{G}_{[t,\infty)} = \sigma\left(\om(A) ; A \in \mathcal{B}([t,\infty) \times \mathbb{R}^d) \right)$ completed by null sets, so
\[\{W_\infty = 0\} \in \bigcap_{t>0}^\infty \mathcal{G}_{[t,\infty)}.\]
The theorem now follows from Komogorov's 0-1 law.
\end{proof}

This dichotomy calls for a definition.

\begin{definition} We say that the polymer is in the \textbf{weak disorder} phase when $W_\infty > 0$ almost surely. We say it is in the \textbf{strong disorder} phase when $W_\infty = 0$ almost surely.
\end{definition}

The phase diagram is connected in the $\beta$-parameter space. 

\begin{theorem}\label{th:313} There exist two critical parameters ${\bar{\b}_c}^- \in [-\infty,0]$ and ${\bar{\b}_c}^+ \in [0,\infty]$, depending only on $\nu,r$ and $d$, such that
\begin{itemize}
\item For all $\b \in ({\bar{\b}_c}^-,{\bar{\b}_c}^+) \,\cup \{0\}$, the polymer belongs to the weak disorder phase.
\item For all $\b \in \mathbb{R}\setminus [{\bar{\b}_c}^-,{\bar{\b}_c}^+]$, the polymer belongs to the strong disorder phase.
\end{itemize}
\end{theorem}
\begin{proof}
Let $\theta$ be a real number in $(0,1)$ and denote $Y_t = W_t^{\theta}$ for all $t\geq 0$. The family $(Y_t)_{t\geq 0}$ is a collection of positive random variables verifying  \[ \sup_{t\geq 0} \IP \big[ Y_t^{1/\theta} \big] = \sup_{t\geq 0} \IP [W_t] = 1 < \infty.\] As $1/\theta$ is strictly greater than $1$, this relation implies the uniform integrablity of $(Y_t)_{t\geq 0}$. Since the process $(W_t^\theta)_{t\geq 0}$ converges almost surely to $W_\infty^\theta$, we get from uniform integrability that
\begin{equation} \label{eq:CV_Wtheta}
\lim_{t\to \infty} \IP \left[W_t^{\theta} \right] = \IP \left[ W_\infty^\theta \right].
\end{equation}
Now, one can observe that the right hand side term is positive if and only if \eqref{eq:WD} holds and that it is zero if and only if \eqref{eq:SD} holds. To prove the theorem, it is then enough to prove that $\b \mapsto \IP \left[ W_\infty^\theta \right]$ is a non-increasing function of $|\b|$ and choose for example 
\begin{equation}
\label{equation:barbetac}
{\bar{\b}_c}^+ = \inf\{\b \geq 0 : \IP \left[ W_\infty^\theta \right] = 0\} \;,
\end{equation}
which does not depend on $\theta \in (0,1)$.
Using \eqref{eq:CV_Wtheta}, we now just have to show that $\b \mapsto \IP \left[ W_t^\theta \right]$ is an non-increasing function of $|\beta|$ for all positive $t$.
By standard arguments, we get that
\begin{align*} \label{eq:derivee_PWt} \frac{\partial}{\partial \b} \IP \left[ W_t^\theta \right] 
& = \IP \left[ \theta W_t^{\theta-1} \frac{\partial}{\partial \b} W_t  \right]\\
& =  \IP \left[ \theta W_t^{\theta-1} P \left[ \left(\om(V_t) - \lambda'(\b) \nu r^d t \right) e^{\b \om(V_t) -\lambda (\b)\nu r^d t}\right] \right]\\
& = \theta \, P \, \IP \left[ W_t^{\theta-1} \left(\om(V_t) - \lambda'(\b) \nu r^d t \right) e^{\b \om(V_t) -\lambda (\b)\nu r^d t} \right].
\end{align*}
Introducing the probability measure $\IP^\b$ on point measures, given by
\[\mathrm d\IP^\b( \om) = e^{\b \om(V_t) - \lambda r^d \nu t} \mathrm d\IP(\om),\]
the derivative of $\IP[ W_t^\theta ]$ is now given by
\begin{equation} \label{eq:derivee_PWt}
\frac{\partial}{\partial \b} \IP \left[ W_t^\theta \right] = \theta \, P \  \IP^\b \left[W_t^{\theta - 1} (\om(V_t) - r^d \nu \lambda' t ) \right].
\end{equation}
In Proposition \ref{prop:PoissonGirsanov} just below, we will see that under the probability measure $\IP^\b$, $\om$ is a Poisson point process on $\mathbb{R}_+ \times \mathbb{R}$.

We can then use the \emph{Harris-FKG inequality} for Poisson processes \cite[th. 11 p. 31]{Last2016} in order to bound the above expectation. Indeed, the variable $\om(V_t) - r^d \nu \lambda' t$ is an increasing function of the point process and by definition, the process $W_t^{\theta - 1}$ is then a decreasing function of $\om$ when $\b \geq 0$ (resp. increasing when $\b < 0$). Applying the FKG inequality, we find that for positive $\b$
\begin{equation} \label{eq:FKG-application}
\IP^\b \left[W_t^{\theta - 1} (\om(V_t) - \lambda' r^d \nu t ) \right ] \leq  \  \IP^\b \left[W_t^{\theta -1} \right] \IP^\b \left[(\om(V_t) - r^d \nu \lambda' t )\right] = 0,
\end{equation}
where the last equality is a result of the relation
\[\IP [\om(V_t) e^{\b \om(V_t)}] =  \lambda'(\b) r^d \nu t.\]
The same result with opposite inequality comes when $\b < 0$.
Thus, we get from \eqref{eq:derivee_PWt} and \eqref{eq:FKG-application} that $\IP[W_t^\theta]$ is a non-increasing function of $|\b|$.
\end{proof}
%%%
We recall at this point that Poisson processes with mutually absolutely continuous intensity measures are themselves mutually absolutely continuous.
%%%
\begin{proposition} \label{prop:PoissonGirsanov} Let $\eta$ be a Poisson point process on a measurable space $E$, of intensity measure $\mu$. Let $f$ be a function such that $e^f-1 \in L^1(\mu)$. Then, under the probability measure $\mathbb{Q}$ defined by 
\[ \frac{\rmd \mathbb{Q}}{\rmd \mathbb{P}} = \exp\left(\int_E f(x) \, \eta(\rmd x) - \int_E (e^{f(x)}-1) \, \rmd \mu(x)\right),\]
the process $\eta$ is a Poisson point process of intensity measure $e^f \, \rmd \mu$.
\end{proposition}
\begin{proof}
Let $g$ be any non-negative measurable function. As the Laplace functional characterizes Poisson processes (theorem 3.9 in \cite{LastPenrose}), we compute it for the point process $\eta$ under the measure $\mathbb{Q}$:
\begin{align*}
\mathbb{Q} \exp\left\{-\int_E g(s) \, \eta(\rmd x)\right\} & = \mathbb{P} \exp \left\{ \int_{E} f(x)-g(x) \, \eta(\rmd x) \right\} e^{ -\int (e^{f(x)}-1) \, \rmd \mu(x) } \\
& =  \exp \left\{ \int_{E} (e^{ f(x)-g(x)}-1) \, \rmd \mu(x) \right\} e^{ -\int (e^{f(x)}-1) \, \rmd \mu(x) }\\
& =  \exp \left\{ \int_{E} (e^{-g(x)}-1) e^{f(x)} \, \rmd \mu(x) \right\},\\
\end{align*}
where the second equality is an application of \eqref{eq:PoissonExpMoment}.
The expression we obtain corresponds, as claimed, to a Poisson point process of intensity measure $e^f \, \rmd \mu$.
\end{proof}


\section{The self-consistency equation and UI properties in the weak disorder}
\subsection{Proof of the self-consistency equation on $W_\infty$} \label{sec:proofSelfconsistent}
In this section, we prove that one can take the limit in the identity $W_{s+t} = P [e_{t} \, W_s \circ \theta_{t,B_t}]$ and obtain the  equation of self-consistency:
\begin{equation}
\label{eq:self_consistent_equation}
W_\infty = P [e_{t} \, W_\infty \circ \theta_{t,B_t} ].
\end{equation}
A part of the problem is that we only have almost sure convergence of the $W_s \circ \theta_{t,x}$ for countable number of $x$'s. 
To deal with this issue, we show that the quantity
\begin{equation}W_t(x) := e^{-\lambda (\b)\nu r^d t} P_x\lef[\exp \lef(\b \om(V_t) \rig) \rig],
\end{equation} 
does not vary too much with $x$, in the sense of the following lemma:


\begin{lemma} \label{lm:distanceWtxWt0}
There exists a constant $C=C(\b,\nu,r)$, such that, for all $t\in [0,\infty]$ and $x,y\in\mathbb{R}^d$,
\begin{equation} \label{eq:distanceWtxWt0}
\IP\left[\left|W_t(x)-W_t(y)\right|\right] \leq C |x-y|.
\end{equation}
\end{lemma}
\begin{proof} To simplify the notations, we only consider the case where $y=0$. To recover the lemma, it is enough to argue that the Poisson environment is invariant in law under a translation in space of vector $y$.

To prove \eqref{eq:distanceWtxWt0} for $y=0$, we write the difference of the two martingales as expectations over two coupled Brownian motions, in the same environment.
The coupling we consider is the mirror coupling, which is defined as follows (see \cite{Hsu2013} for more details).

At time $0$, one of the Brownian motions is starting at $0$ and one is starting at $x \in \mathbb{R}^d$. Denote by $H$ be the hyperplane bisecting the segment $[0,x]$, which is the hyperplane passing by $x/2$ and orthogonal to the vector $x$. Let also
\[\tau = \inf\{t\geq 0 | B_t \in H\},\]
be the first hitting time of $H$ by $B$. Then, define $\widetilde{B}$ as the path that coincides with the reflection of path of $B$ with respect to $H$ for times before $\tau$, and that coincides with $B$ after $\tau$. 

%%%%%%
\begin{figure} \centering
\includegraphics[scale=0.45]{mirror-coupling.eps}
\vskip -3cm
\caption{Mirror coupling in $d=1$.} 
 \label{fig:mirror}
\end{figure}
%%%%%%%%
%\bfc [ quick drawing ???] \ec
The process $\widetilde B$ has the law of a Brownian motion starting from $x$. Moreover, the time $\tau$ is the first time $B$ and $\widetilde B$ meet. After $\tau$, the processes coincide. The variable $\tau$ has the following cumulative distribution function:
\begin{equation} \label{eq:cumulativeOfTau}
P(\tau \geq z)=\phi_z(|x|),
\end{equation}
where, for positive $z$, \[\phi_z(|x|)=\frac{2}{\sqrt{2\pi z}}\int_{0}^{|x|/2} e^ {-u^2/2z}\rmd u,\] and where $|\cdot|$ is the Euclidean distance. In the litterature, a coupling that satisfies this relation is said to be maximal (see \cite{Hsu2013}).
%Let $e_t = \exp(\b \om(V_t(B)) - \lambda(\b)\nu r^d t)$ and let \linebreak $\tilde{e}_t = \exp(\b \om(V_t(\widetilde B)) - \lambda(\b)\nu r^d t)$, which 
%we factorize in the contributions before and after $B$ and $\widetilde B$ coalesce, so that, for $t\in [0,\infty)$,
Let $$e_t = \exp(\b \om(V_t(B)) - \lambda(\b)\nu r^d t), \quad \tilde{e}_t = \exp(\b \om(V_t(\widetilde B)) - \lambda(\b)\nu r^d t),$$ which we factorize in the contributions before and after $B$ and $\widetilde B$ coalesce, so that, for $t\in [0,\infty)$,
\begin{align*}
\IP[|W_t(x)-W_t(0)|] &\leq \IP P[|\tilde{e}_t - e_t|] \\
& = P\IP [|(\tilde{e}_{t\wedge \tau}-{e}_{t\wedge \tau})e_{(t-t\wedge \tau)^+}\circ \theta_{t\wedge \tau,B_{t\wedge \tau}}|]\\
& = P\IP [|\tilde{e}_{t\wedge \tau}-{e}_{t\wedge \tau}|],
\end{align*}
where the last equality is a result of the independance of the Poisson environment before and strictly after time $t\wedge \tau$. 

Then, we distinguish the cases where $B$ or $\widetilde{B}$ encounter a point of the environment before $t\wedge \tau$, and the cases where they don't. We get that $P\IP [|\tilde{e}_{t\wedge \tau}-{e}_{t\wedge \tau}|]$ writes:
\begin{align*}
& P\IP \left[|e_{t\wedge \tau} - \tilde{e}_ {t\wedge \tau}| \mathbf{1}_{\{\om(V_{t\wedge \tau}(B))>1,\,\om(V_{t\wedge \tau}( \widetilde B))>1\}}\right] + P\IP\left[(e_{t\wedge \tau} - e^{-\lambda \nu r^d t\wedge \tau}) \mathbf{1}_{\{\om(V_{t\wedge \tau}(B))>1,\,\om(V_{t\wedge \tau}( \widetilde B))=0\}}\right]\\
& \qquad + P\IP\left[(\tilde{e}_{t\wedge \tau} - e^{-\lambda \nu r^d t\wedge \tau}) \mathbf{1}_{\{\om(V_{t\wedge \tau}(B)=0,\,\om(V_{t\wedge \tau}( \widetilde B)>1\}}\right].
\end{align*}
We first use the triangle inequality in the first expectation of the sum, and neglect the negative terms in two other expectations. Then, recombining the terms, one obtains that
%\begin{align}
%\IP[|W_t(x)-W_t(0)|] & \leq 2 P\IP \left[e_{t\wedge \tau} \mathbf{1}_{\{\om(V_{t\wedge \tau}(B))>1\}}\right] + 2 P\IP \left[\tilde{e}_{t\wedge \tau} \mathbf{1}_{\{\om(V_{t\wedge \tau}(\widetilde{B}))>1\}}\right] \nonumber \\ 
%& = 4 P\IP \left[e_{t\wedge \tau} \mathbf{1}_{\{\om(V_{t\wedge \tau}(B))>1\}}\right] \nonumber \\
%\label{eq:ineqInLemmaWtx} & \leq 4 P\IP \left[e_{\tau} \mathbf{1}_{\{\om(V_{\tau}(B))>1\}}\right],
%\end{align}

\begin{align}
\IP[|W_t(x)-W_t(0)|] & \leq 2 P\IP \left[e_{t\wedge \tau} \mathbf{1}_{\{\om(V_{t\wedge \tau}(B))>1\}}\right] + 2 P\IP \left[\tilde{e}_{t\wedge \tau} \mathbf{1}_{\{\om(V_{t\wedge \tau}(\widetilde{B}))>1\}}\right] \nonumber \\ 
& = 4 P\IP \left[e_{t\wedge \tau} \mathbf{1}_{\{\om(V_{t\wedge \tau}(B))>1\}}\right]  \label{eq:ineqInLemmaWtx}
\end{align}
where the equality is a consequence of invariance in law of the Poisson environment under the translation by $x$.

For any $t$, the variable $\om(V_t)$ is a Poisson r.v. of parameter $ \nu r^d t$, so that
\begin{equation*}
P\IP \left[e_{t\wedge\tau} \mathbf{1}_{\{\om(V_{t\wedge\tau}(B)>1\}}\right]  = \sum_{k=1}^\infty P \IP [e_{t\wedge\tau} \mathbf{1}_{\{\om(V_{t\wedge\tau})=k\}}]
 = \sum_{k=1}^\infty P \left[e^{\b k-\lambda\nu r^d (t\wedge\tau)} \frac{(\nu r^d (t\wedge\tau))^k}{k!}e^{-\nu r^d (t\wedge\tau)} \right],
\end{equation*}
from which we get by standard computations that
\begin{equation} \label{eq:expecInLemmaWtx}
P\IP \left[e_{t\wedge\tau} \mathbf{1}_{\{\om(V_{(t\wedge\tau)}(B)>1\}}\right] = P[1-\exp(-e^\b \nu r^d t\wedge\tau)] \leq P[1-\exp(-e^\b \nu r^d \tau)].\end{equation}

To control this last expectation, we will first notice that $\tau / |x|^2$ is independent of $|x|$, and we will compute its density. By equation \eqref{eq:cumulativeOfTau} and the change of variable $u=|x|\sqrt(z) v$, we get that
\begin{align*}
\P(\tau/|x|^2 > z) = \frac{1}{\sqrt{2\pi}} \int_0^{1/2\sqrt{z}} e^{-v^2/2} \rmd v.
\end{align*}
This function of $z$ is continuous and everywhere differentiable on $[0, \8)$. %, except maybe at $z=0$ [??? check this]. 
The variable $\tau/|x|^2$ hence admits a density with respect to the Lebesgue measure, given by
\[f(z) = \frac{1}{4\sqrt{2\pi}} \frac{e^{-1/8z}}{z^{3/2}}.\]
Therefore, one gets that the right-hand side of \eqref{eq:expecInLemmaWtx} can be written as
\begin{align*}
&P\left[\left(1-\exp(-e^\b \nu r^d {\tau}{|x|^{-2}} |x|^2)\right)\mathbf{1}_{\{\tau/|x|^2 \leq 1\}}\right] + P\left[\left(1-\exp(-e^\b \nu r^d {\tau}{|x|^{-2}} |x|^2)\right)\mathbf{1}_{\{\tau/|x|^2 > 1\}}\right]\\
&\qquad \qquad \qquad \qquad 
\leq e^\b \nu r^d |x|^2 + \int_1^\infty \left(1-\exp(-e^\b \nu r^d z |x|^2)\right) f(z) \rmd z.
\end{align*}
As there is some constant $C>0$ such that $f(z) \leq C z^{-3/2}$, after the change of variables $z|x|^2=u$, the integral above can be bounded by
\[|x| \int_{0}^\infty C \left(1-\exp\left(-e^\b \nu r^d u \right)\right) v^{-3/2} \rmd v,\]
where one can check that the integral converges. Since its value only depends on $\b,\nu$ and $r$, we finally get that there exists some constant $C'=C(\b,\nu,r)$, such that
\[\IP[|W_t(x)-W_t(0)|] \leq 4P[1-\exp(-e^\b \nu r^d \tau)] \leq C' (|x|^2 + |x|),\]
where the first inequality comes from combining \eqref{eq:ineqInLemmaWtx} and \eqref{eq:expecInLemmaWtx}. Since the $W_t(x)$ variables have bounded expectations, this proves the lemma in the case where $t<\infty$. The $t=\infty$ case is then a consequence of Fatou's lemma.
\end{proof}

\vspace{0.4cm}
We can now show that the self-consistency equation holds. Let $\delta >0$ be a parameter that will go to $0$. For all $q \in \mathbb{Z}^d$, define $\Delta( q)$ to be the cube of length $\delta$ centered at $\delta q$, so that all the cubes form a partition of the space $\mathbb{R}^d$.
We get that the right-hand side of \eqref{eq:WtIntegralEquation} satisfies
\begin{align}
P [e_{t} \, W_s \circ \theta_{t,B_t} ] & = \sum_{q\in\mathbb{Z}^d} P [e_{t} \, W_s \circ \theta_{t,B_t} ; B_t\in \Delta( q) ] \nonumber \\
& = \sum_{q\in\mathbb{Z}^d} P [e_{t} \, W_s \circ \theta_{t,\delta q} ; B_t\in \Delta( q) ] + A_s^\delta, \label{eq:approxInL1normForWSD}
\end{align}
where
$A_s^\delta = \sum_{q\in\mathbb{Z}^d} P [e_{t} \, (W_s \circ \theta_{t,B_t}-W_s \circ \theta_{t,\delta q}) ; B_t\in \Delta( q) ].$

First observe that $\IP$-almost surely, $W_s \circ \theta_{t,{\delta q}}$ converges  to $W_\infty \circ \theta_{t,{\delta q}}$ for all $q\in\mathbf{Z}^d$, so Fatou's lemma entails
\[\IP\text{-a.s.}, \quad \liminf_{s\to\infty} \sum_{q\in\mathbb{Z}^d} P [e_{t} \, W_s \circ \theta_{t,\delta q} ; B_t\in \Delta( q) ] \geq \sum_{q\in\mathbb{Z}^d} P [e_{t} \, W_\infty \circ \theta_{t,\delta q} ; B_t\in \Delta( q) ],\] so that, by \eqref{eq:approxInL1normForWSD} and letting $s\to\infty$ in \eqref{eq:WtIntegralEquation},
\begin{equation} \label{eq:ineqOnWinftyAfterFatou}
\IP\text{-a.s.}, \quad W_\infty \geq \sum_{q\in\mathbb{Z}^d} P [e_{t} \, W_\infty \circ \theta_{t,\delta q} ; B_t\in \Delta( q) ] + \liminf_{s\to\infty} A_s^\delta.
\end{equation}

Furthermore, using the fact that $W$ is a martingale, one can check that $s\mapsto A_s^\delta$ is also a martingale with respect to the filtration $\{\mathcal{G}_{t+s}\}_{s\geq 0}$.
% [is it ok???]. 
For any time $S\geq 0$, Lemma \ref{lm:distanceWtxWt0} implies that it satisfies
\begin{align*}
\IP[|A_S^\delta|] & \leq \sum_{q\in\mathbb{Z}^d} P \left[\IP[e_{t}\,|W_S \circ \theta_{t,B_t}- W_S \circ \theta_{t,\delta q}|] \mathbf{1}_{B_t\in \Delta( q)} \right]\\
&\leq C(\delta^2 + \delta) \sum_{q\in\mathbb{Z}^d} P\left[ \mathbf{1}_{B_t\in \Delta( q)}\right]\\ &= C(\delta^2 + \delta),
\end{align*}
where, in the second inequality, we have factorized by $\IP[e_{t}] = 1$, using the independence under $\IP$ of the environment before and strictly after time $t$. Thus, by Doob's inequality \cite[Th. 3.8 (i), Ch.~1]{KS91}, we have for all $u>0$,
\[\IP\left[\sup_{0\leq s\leq S} |A_s^\delta| \geq u \right] \leq \frac{C(\delta^2 + \delta)}{u},\]
where we can let $S\to\infty$ by monotone convergence. This implies that $\sup_{s\geq 0} |A_s^\delta|$ converges in probability to $0$, when $\delta \to 0$, which in turn implies that
\[\liminf_{s\to\infty} |A_s^\delta| \overset{\IP}{\longrightarrow} 0.\]

Then, using Lemma \ref{lm:distanceWtxWt0} in the case where $t=\infty$, the same computation as above would show that
\[\IP\left[ \left| \sum_{q\in\mathbb{Z}^d}P \left[e_{t} \, W_\infty \circ \theta_{t,B_t} ; B_t\in \Delta( q)\right]- \sum_{q\in\mathbb{Z}^d} P \left[e_t \, W_\infty \circ \theta_{t,\delta q} ; B_t\in \Delta( q)\right] \right| \right] \leq C(\delta^2 + \delta),\]
and hence,
\[\sum_{q\in\mathbb{Z}^d} P \left[e_t \, W_\infty \circ \theta_{t,\delta q} ; B_t\in \Delta( q)\right]   \overset{L^1}{\longrightarrow} P [e_{t} \, W_\infty \circ \theta_{t,B_t}].\]
In particular, we have shown that the right-hand side of \eqref{eq:ineqOnWinftyAfterFatou} converges in probability to $P [e_{t} \, W_\infty \circ \theta_{t,B_t}]$, so that almost surely
\[ W_\infty \geq  P [e_{t} \, W_\infty \circ \theta_{t,B_t}].\]
Observe that the two quantities have the same expectations to conclude that \eqref{eq:self_consistent_equation} holds.


\subsection{Uniform integrability in the weak disorder}
\begin{proposition} \label{prop:UI}
The martingale $W_t$ is uniformly integrable if and only if the polymer is in the weak disorder phase, i.e. $W_\infty > 0,\, \IP$-almost surely.
\end{proposition}
\begin{proof}
If $W_t$ is UI, then $W_t$ converges in $L^1$ to $W_\infty$, so that $\IP[W_\infty] = 1 > 0$, therefore weak disorder must hold by the dichotomy.

Suppose now that the polymer is in weak disorder and set
\[X_{t,x} = \frac{W_\infty\circ \theta_{t,x}}{\IP [W_\infty]}.\]
The self-consistency equation \eqref{eq:self_consistent_equation} writes
$X_{0,0} = P[e_t X_{t,B_t}]$,
so that, as $X_{t,x}$ is independent of $\mathcal{G}_t$ for all $x$,
\[\IP[ X_{0,0} | \mathcal{G}_t] = P\left[e_t \IP[X_{t,B_t}]\right] = P[e_t].\]
This shows that for all $t\geq 0$, 
\[W_t = \IP[X_{0,0} | \mathcal{G}_t] \quad \text{a.s.}\]
Hence, $(W_t)_{t\geq 0}$ is uniformly integrable since the family of the right-hand side is a uniformly integrable martingale.
\end{proof}
















\section{The $\bf{L^2}$-region}
\begin{theorem} \label{th:LTwoRegion} (i) There exist two critical parameters ${\b_2}^- \in [-\infty,0]$ and ${\b_2}^+ \in [0,\infty]$, depending only on $\nu,r$ and $d$, such that, if $\b \in ({\b_2}^-,{\b_2}^+) \, \cup \{0\}$ then
\begin{equation}\label{eq:supWt2} \sup_{t\in\mathbb{R}} \IP [W_t^2] < \infty,
\end{equation} and such that the supremum is infinite if $\b \in \mathbb{R}\setminus [{\b}_2^-,\b_2^+]$.

(ii) Furthermore, if $d\geq 3$, there exists a constant $c(d)\in (0,\infty)$, such that \eqref{eq:supWt2} holds whenever
\begin{equation} \label{eq:L2_condition}
\lambda(\b)^2 \nu r^{d+2} < c(d),
\end{equation}

(iii) In particular, ${\b_2}^- < 0$ and ${\b_2}^+ > 0$ whenever $d\geq 3$.

(iv) Also for  $d\geq 3$, when $\nu r^{d+2} < c(d)$ the constant in \eqref{eq:L2_condition}, we have ${\b_2}^- =-\8$.
\end{theorem}

\begin{definition}We call the $\,\textbf{$\bf{L^2}$-region}$ the set of parameters $\beta, \nu, r$ for which \eqref{eq:supWt2} holds.\end{definition}

\begin{proof}[Proof of Theorem \ref{th:LTwoRegion}]
We introduce the product measure $P^{\otimes 2}$ of two independent Brownian motions $B_t$ and $\tilde{B}_t$ starting from $0$ with respective tubes $V_t$ and $\tilde{V}_t$. The main idea is to write
\[W_t^2 = P^{\otimes 2} [e^{\b \om(V_t)} e^{\b \om(\tilde{V}_t)}]e^{-2\lambda \nu r^d t},\]
so that, using Fubini's theorem,
\begin{equation} \IP[W_t^2] = P^{\otimes 2} \, \IP [e^{\b (\om(V_t)+ \om(\tilde{V}_t)}]e^{-2\lambda \nu r^d t}.
\end{equation}
One can see that \[\om(V_t) + \om(\tilde{V}_t) = 2\om(V_t \cap \tilde{V}_t) + \om(V_t \Delta \tilde{V}t),\]
which is the sum of two independent Poisson random variables; computing their Laplace transforms leads us to
%\begin{align}
%&\IP [e^{ 2 \b \om(V_t \cap \tilde{V}_t) + \om(V_t \Delta \tilde{V}t)}]e^{-2\lambda \nu r^d t} \nonumber \\
%&= \exp \left( \lambda(2\b) \nu |V_t \cap \tilde{V}_t| + \lambda \nu |V_t \Delta \tilde{V}t| - 2\lambda \nu r^d t \right) \nonumber \\
%&= \exp\left(\lambda^2 \nu |V_t \cap \tilde{V}_t| \right) \nonumber \\
%\label{eq:cvWt2} &\underset{t\to\infty}{\longrightarrow} \exp\left(\lambda^2 \nu |V_\infty \cap \tilde{V}_\infty| \right), 
%\end{align}
\begin{eqnarray}
%\IP [e^{ 2 \b \om(V_t \cap \tilde{V}_t) + \om(V_t \Delta \tilde{V}t)}]e^{-2\lambda \nu r^d t} 
 \IP[W_t^2]
\nonumber 
&=& \exp \left( \lambda(2\b) \nu |V_t \cap \tilde{V}_t| + \lambda \nu |V_t \Delta \tilde{V}_t| - 2\lambda \nu r^d t \right) \nonumber \\
&=& \exp\left(\lambda^2 \nu |V_t \cap \tilde{V}_t| \right) \nonumber \\
\label{eq:cvWt2} &\underset{t\to\infty}{\longrightarrow}& \exp\left(\lambda^2 \nu |V_\infty \cap \tilde{V}_\infty| \right), 
\end{eqnarray}
where the second equality is obtained using $|V_t|=|\tilde{V}_t|=tr^d$ and $\lambda(\b)^2=\lambda(2\b)-2\lambda(\b)$, while the limit is justified by monotone convergence.
Now,
\begin{align*} |V_\infty\cap \tilde{V}_\infty| &= \int_0^\infty |U(B_t)\cap U(\tilde{B}_t)|\rmd t \\
&= \int_0^\infty |U(0)\cap U(\tilde{B}_t-B_t)|\rmd t\\
& \eqlaw \int_0^\infty |U(0)\cap U(B_{2t})|\rmd t,
\end{align*}
since $(B_t + \tilde{B}_t)_{t\geq 0} \eqlaw (B_{2t})_{t\geq 0}$. Hence, using monotone convergence and \eqref{eq:cvWt2}, we get that
\begin{equation} \label{eq:supESPWt2}
\sup_{t\in\mathbb{R}} \IP [W_t^2] = \lim_{t\to\infty} \IP [W_t^2] = P\left[\exp\left(\frac{\lambda(\b)^2}{2} \nu \int_0^\infty |U(0)\cap U(B_{t})|  \rmd t\right) \right],
\end{equation}
where the first equality is a consequence of $(W_t^2)_{t\geq 0}$ being a submartingale. Equation \eqref{eq:supESPWt2} shows that $\sup_{t} \IP[W_t^2]$ is an increasing function of $|\b|$, which proves part (i) of the theorem.

To prove the second part, we will bound the right hand side of \eqref{eq:supESPWt2}. First observe that \[|U(0)\cap U(B_t)| \leq |U(0)| \, \mathbf{1}_{|B_t| \leq 2\gamma_d^{-1/d}r} \eqlaw r^d \mathbf{1} \left\{\left|B\left(\frac{t\gamma_d^{2/d}}{4r^{2}}\right)\right| \leq 1 \right\},\]
so that, by a change of variables,
\begin{equation} \label{eq:bound_supWt2} \sup_{t\in\mathbb{R}} \IP [W_t^2] \leq P\left[\exp\left(2\lambda(\b)^2 \nu r^{2+d}\gamma_d^{-2/d} \int_0^\infty \mathbf{1}_{|B_t| \leq 1} \, \rmd t\right) \right].
\end{equation}
When $d\geq 3$, the Brownian motion is transient and has the following property:
\[\alpha(d) = \sup_{x\in \mathbb{R}^d} P_x \left[ \int_0^\infty \mathbf{1}_{|B_t| \leq 1}\,\rmd t \right] < \infty.\]
By Khas'minskii's lemma \cite[p. 8, Lemma 2.1]{Sznitman}, this implies that
\[P \left[ \exp\left(u \int_0^\infty \mathbf{1}_{|B_t| \leq 1} \, \rmd t\right) \right] < (1-u\alpha)^{-1},\]
whenever $u\alpha < 1$. Looking back at \eqref{eq:bound_supWt2}, this condition finally leads to \eqref{eq:L2_condition}.

Part (iii) is obtained by observing that $\lambda(\beta) \to 0$ as $\beta \to 0$, so that condition \eqref{eq:L2_condition} is fulfilled for small enough $\b$.

To prove (iv), just note that the right-hand side of  \eqref{eq:L2_condition} tends to $\nu r^{d+2}$ as $\b \to -\8$.
\end{proof}

\section{Relations between the different critical temperatures}

The critical values $\b_2^\pm, {\bar{\b}_c}^\pm$ and $  \b_c^\pm$ defined in Theorems \ref{defdef..}, \ref{th:313} and
\ref{th:LTwoRegion}
are ordered. 

\begin{proposition} The following properties hold:
\begin{enumerate}[label=(\roman*)]
\item For all $d\geq 1$,
\begin{equation}
\b_2^+ \leq {\bar{\b}_c}^+ \leq \b_c^+ < \infty \quad \text{and} \quad \b_c^- \leq {\bar{\b}_c}^- \leq \b_2^-.
\end{equation}
\item When $d\geq 3$, these parameters are all non-zero.
\item When $d\geq 3$ and $\nu r^{d+2}$ is small enough, $\b_c^- = {\bar{\b}_c}^- = \b_2^- = -\infty$.
\end{enumerate} 
\end{proposition}
\begin{remark} Point (iii) tells us that if the intensity of the Poisson point process or the radius are sufficiently small, the polymer will not really be impacted by the environment.
\end{remark}
\begin{remark}\label{rk:criticd=12}
For $d=1,2$, it holds that
%[references??? Lacoin, Bertin] \cite[Th. 2.1.1 (b)]{CY05}
that $\b_2^\pm = {\bar{\b}_c}^\pm = \b_c^\pm = 0$. The reader is referred to the proofs in \cite{BeLa15, Bertin, CY05, Lacoin10} for other similar models, and can be convinced that the arguments go through in our case \cite{BertinUnp}.
\end{remark}
%\begin{remark}
%We have $\b_c^+ < \infty$ for all $d \geq 1$.  
%\bfc [REFE...] \ec
%However, $\b_c^-$, is equal to $-\8$ in dimension $d \geq 3$ for some $\nu$'s.
%\end{remark}
\begin{remark}
A long-standing conjecture  is that $\b_c^\pm = {\bar{\b}_c}^\pm$, i.e., that the weak/strong disorder transition coincide with the high/low temperature one. % (except possibly on the critical line). 
\end{remark}


\begin{proof}
We first show that $\b_2^+ \leq {\bar{\b}_c}^+$.
Suppose $\b_2^+>0$ and let $0\leq \b < \b_2^+$. By definition, we have
\[\sup_{t\geq 0} \IP [W_t^2] < \infty,\]
so that $(W_t)_{t\geq 0}$ is a martingale bounded in $L^2$. Thus, $(W_t)_{t\geq 0}$ converges in $L^2$ norm, which implies $L^1$ convergence.

Since $\IP [W_t] = 1$ for all $t$, we get that $\IP [W_\infty] = 1$, so \eqref{eq:WD} must hold and hence $\b \leq {\bar{\b}_c}^+$. As it is true for all $\b<\b_2^+$, the desired inequality follows directly.

We now turn to the proof of ${\bar{\b}_c}^+ \leq \b_c^+$. Again, suppose that ${\bar{\b}_c}^+>0$ and let $0\leq \b < {\bar{\b}_c}^+$. We have $W_\infty >0$ almost surely, so $\ln W_t \to \ln W_\infty$ almost surely, that is to say
\[\ln Z_t - \lambda \nu r^d t \underset{t\to\infty}{\longrightarrow} \ln W_\infty.\]
Dividing by $t$, we get that
\[\frac{1}{t}\ln Z_t - \lambda \nu r^d \underset{t\to\infty}{\longrightarrow} 0,\]
thus $p(\b, \nu, r) = \lambda \nu r^d$, i.e. $\b \leq \b_c$.
The same argument goes for the negative critical values. This ends the proof of (i).

Points (ii) and (iii) are repeated from Theorem \ref{th:LTwoRegion}. 
\end{proof}




 

 